\documentclass[12pt]{article}
\usepackage[utf8]{inputenc}
\usepackage[T1]{fontenc}
\usepackage[margin=1in]{geometry}
\usepackage{setspace}
\usepackage{xcolor}
\usepackage[colorlinks=true, urlcolor=blue!60!black, linkcolor=blue!60!black]{hyperref}
\setlength{\parindent}{2em}
\setlength{\parskip}{1.25em}
\renewcommand{\baselinestretch}{1.0}
\begin{document}
\begin{center}
{\LARGE \textit{Wisconsin v. Yoder}, 406 U.S. 205 (1972)\par}
\vspace{0.8em}
{\large Case Brief\par}
\end{center}
\vspace{1em}

\section{Memory Jogger}
Applying strict scrutiny under the Free Exercise Clause, the Court held that Wisconsin could not criminally punish Amish parents for refusing to send their children to formal high school after the eighth grade, because the State's interest in one or two more years of schooling did not outweigh the grave threat to the Amish religious way of life (\href{https://supreme.justia.com/cases/federal/us/406/205/}{Justia}).

\section{Detailed Case Facts}
Respondents Jonas Yoder and Wallace Miller were members of the Old Order Amish religion, and respondent Adin Yutzy was a member of the Conservative Amish Mennonite Church; all three families resided in Green County, Wisconsin (\href{https://supreme.justia.com/cases/federal/us/406/205/}{406 U.S. at 207}). Wisconsin's compulsory school attendance statute, Wis. Stat. \S~118.15 (1969), required parents to cause their children between ages 7 and 16 to attend public or private school regularly until the end of the school term in which the child turned 16, subject to fines of \$5 to \$50, imprisonment for up to three months, or both (\href{https://supreme.justia.com/cases/federal/us/406/205/}{406 U.S. at 207 n.2}).

Respondents declined to enroll their children---Frieda Yoder (age 15), Barbara Miller (age 15), and Vernon Yutzy (age 14)---in public or private high school after the children graduated from the eighth grade of public school (\href{https://supreme.justia.com/cases/federal/us/406/205/}{406 U.S. at 207 \& n.1}). Before trial, respondents' counsel proposed a compromise to the State Superintendent of Public Instruction modeled on the ``Pennsylvania plan,'' under which Amish adolescents attended an Amish vocational school three hours per week for English, mathematics, health, and social studies while completing supervised agricultural and homemaking projects on the family farm (\href{https://supreme.justia.com/cases/federal/us/406/205/}{406 U.S. at 208 n.3}). The Superintendent rejected the proposal as failing to provide ``substantially equivalent education'' (\href{https://supreme.justia.com/cases/federal/us/406/205/}{406 U.S. at 208 n.3}).

At trial, the State stipulated that respondents' religious beliefs were sincere (\href{https://supreme.justia.com/cases/federal/us/406/205/}{406 U.S. at 209}). Uncontradicted expert testimony by Dr. John Hostetler (a sociologist of Amish society) and Dr. Donald A. Erickson (an education scholar) established that Old Order Amish faith---rooted in sixteenth-century Swiss Anabaptism and a literal reading of Paul's Epistle to the Romans, ``be not conformed to this world''---pervades every aspect of daily life through the community's \emph{Ordnung} (\href{https://supreme.justia.com/cases/federal/us/406/205/}{406 U.S. at 209--12, 216}). While the Amish accepted formal elementary schooling through the eighth grade to acquire basic literacy and arithmetic (``the three R's''), they opposed formal high school because it removed adolescents physically and emotionally from the community during the crucial formative years preceding adult baptism, exposed them to competitive, worldly values, and threatened both individual salvation and the survival of the Amish community (\href{https://supreme.justia.com/cases/federal/us/406/205/}{406 U.S. at 210--12}). Dr. Erickson testified that the Amish system of informal vocational ``learning through doing'' was ``ideal'' for preparing youth for productive, self-sufficient community life (\href{https://supreme.justia.com/cases/federal/us/406/205/}{406 U.S. at 212, 223}). Only one child, Frieda Yoder, testified at trial; she confirmed that her own religious beliefs guided her decision not to attend high school (\href{https://supreme.justia.com/cases/federal/us/406/205/}{406 U.S. at 231 n.21, 237}).

\section{Procedural History}
On complaint of the public school district administrator, respondents were charged, tried, and convicted in Green County Court of violating Wis. Stat. \S~118.15 and fined \$5 each (\href{https://supreme.justia.com/cases/federal/us/406/205/}{406 U.S. at 208}). Although the trial court found that the statute interfered with respondents' sincere religious beliefs, it upheld the law as a ``reasonable and constitutional'' exercise of state power and denied their motion to dismiss (\href{https://supreme.justia.com/cases/federal/us/406/205/}{406 U.S. at 213}). The Wisconsin Circuit Court affirmed the convictions (\href{https://supreme.justia.com/cases/federal/us/406/205/}{406 U.S. at 213}). The Wisconsin Supreme Court reversed, holding that Wisconsin had failed to demonstrate that its interest in establishing and maintaining an educational system overrode respondents' First Amendment Free Exercise rights, 49 Wis. 2d 430, 447, 182 N.W.2d 539, 547 (1971) (\href{https://supreme.justia.com/cases/federal/us/406/205/}{406 U.S. at 213}). The U.S. Supreme Court granted Wisconsin's petition for certiorari and affirmed (\href{https://supreme.justia.com/cases/federal/us/406/205/}{406 U.S. at 207, 236}).

\section{Judicial Votes}
The Supreme Court \textbf{affirmed} the Wisconsin Supreme Court (reversing respondents' convictions) by a \textbf{6--1} vote, with seven Justices participating. The case was argued on December 8, 1971, and decided on May 15, 1972 (\href{https://supreme.justia.com/cases/federal/us/406/205/}{Justia}; \href{https://www.law.cornell.edu/supremecourt/text/406/205}{Cornell LII}):

\begin{itemize}
    \item \textbf{Majority Opinion:} Chief Justice Burger delivered the opinion of the Court, joined by Justices Brennan, Stewart, White, Marshall, and Blackmun.
    \item \textbf{Concurring Opinions:} Justice Stewart filed a concurring opinion, joined by Justice Brennan. Justice White filed a concurring opinion, joined by Justices Brennan and Stewart.
    \item \textbf{Dissenting Opinion:} Justice Douglas filed an opinion dissenting in part (concurring in the judgment as to Jonas Yoder and dissenting as to Adin Yutzy and Wallace Miller).
    \item \textbf{Not Participating:} Justices Powell and Rehnquist took no part in the consideration or decision of the case.
\end{itemize}

\section{Holding}
\textbf{Questions Presented:}

\begin{enumerate}
  \item Does the Free Exercise Clause of the First Amendment, as incorporated against the States by the Fourteenth Amendment, prohibit a State from criminally prosecuting Old Order Amish and Conservative Amish Mennonite parents under a facially neutral, generally applicable compulsory school attendance law for refusing to send their children to formal high school past the eighth grade, where compliance would gravely endanger their religious faith and community survival?
  \item May a State invoke its \emph{parens patriae} authority over children's welfare and secondary education to override parents' religiously grounded objections where no physical or mental harm to the child and no parent-child conflict appears in the record?
\end{enumerate}

\textbf{Holding:} \textbf{Yes} as to Issue~1; \textbf{No} as to Issue~2. Applying strict scrutiny, the Supreme Court held that the First and Fourteenth Amendments prevent Wisconsin from compelling Amish parents to cause their children to attend formal high school until age 16 (\href{https://supreme.justia.com/cases/federal/us/406/205/}{406 U.S. at 234}). Only state ``interests of the highest order and those not otherwise served'' can overbalance legitimate Free Exercise claims (\href{https://supreme.justia.com/cases/federal/us/406/205/}{406 U.S. at 215}). Given the severity of the burden on a centuries-old religious way of life, the Amish acceptance of eight years of formal elementary schooling plus an effective informal vocational alternative, and the absence of demonstrated harm to the children or society, Wisconsin failed to show with particularity how granting an exemption for one or two additional years would adversely affect its educational objectives (\href{https://supreme.justia.com/cases/federal/us/406/205/}{406 U.S. at 235--36}).

\section{Analysis of Opinions}
\subsection{Majority Opinion (Chief Justice Burger)}

Chief Justice Burger's opinion for the Court proceeded in five structured parts:

\subsection{Part I: Balancing Universal Education Against Free Exercise and Parental Rights}

While acknowledging that ``[p]roviding public schools ranks at the very apex of the function of a State,'' the Court emphasized that this interest is not immune from constitutional balancing when it impinges on fundamental rights (\href{https://supreme.justia.com/cases/federal/us/406/205/}{406 U.S. at 213--14}). Drawing on \href{https://supreme.justia.com/cases/federal/us/268/510/}{\emph{Pierce v. Society of Sisters}}, 268 U.S. 510, 534--35 (1925), the Court linked the Free Exercise Clause with the ``traditional interest of parents with respect to the religious upbringing of their children'' (\href{https://supreme.justia.com/cases/federal/us/406/205/}{406 U.S. at 214}). Relying on \href{https://supreme.justia.com/cases/federal/us/374/398/}{\emph{Sherbert v. Verner}}, 374 U.S. 398 (1963), the Court announced its demanding standard: ``only those interests of the highest order and those not otherwise served can overbalance legitimate claims to the free exercise of religion'' (\href{https://supreme.justia.com/cases/federal/us/406/205/}{406 U.S. at 215}).

\subsection{Part II: Religious vs.\ Philosophical Claims and Severity of the Burden}

The Court drew a sharp threshold boundary between religious and secular claims. A ``virtuous and admirable'' way of life based on purely secular or philosophical convictions---analogous to Henry David Thoreau's retreat to Walden Pond---``does not rise to the demands of the Religion Clauses'' (\href{https://supreme.justia.com/cases/federal/us/406/205/}{406 U.S. at 215--16}). Here, however, the Amish way of life was not a matter of personal preference, but ``one of deep religious conviction, shared by an organized group, and intimately related to daily living'' for nearly three centuries (\href{https://supreme.justia.com/cases/federal/us/406/205/}{406 U.S. at 216, 219}). Because consolidated modern high schools exposed Amish youth to worldly values and disrupted adolescent integration into the faith community, Wisconsin's criminal mandate posed an ``objective danger'' and ``a very real threat of undermining the Amish community and religious practice as they exist today,'' forcing the Amish either to abandon their faith or to migrate (\href{https://supreme.justia.com/cases/federal/us/406/205/}{406 U.S. at 218}).

\subsection{Part III: Rejection of the Belief--Action Dichotomy and Facial-Neutrality Defense; Scrutiny of State Interests}

Wisconsin advanced two categorical defenses: (1) that the Free Exercise Clause protects beliefs but not religiously motivated ``actions,'' and (2) that a facially neutral, uniformly applied statute motivated by secular goals cannot violate the Clause. The Court rejected both. First, ``in this context, belief and action cannot be neatly confined in logic-tight compartments'' (\href{https://supreme.justia.com/cases/federal/us/406/205/}{406 U.S. at 220}). Second, ``[a] regulation neutral on its face may, in its application, nonetheless offend the constitutional requirement for governmental neutrality if it unduly burdens the free exercise of religion'' (\href{https://supreme.justia.com/cases/federal/us/406/205/}{406 U.S. at 220}, citing \href{https://supreme.justia.com/cases/federal/us/374/398/}{\emph{Sherbert}}, 374 U.S. 398).

Turning to the State's asserted compelling interests---preparing citizens for intelligent democratic participation and fostering economic self-sufficiency---the Court refused to evaluate those goals in the abstract, insisting instead on examining the marginal impediment caused by an Amish exemption (\href{https://supreme.justia.com/cases/federal/us/406/205/}{406 U.S. at 221}). The record refuted Wisconsin's charge that the Amish fostered ``ignorance'': the Amish were law-abiding, self-sufficient, rejected public welfare, and had even been exempted by Congress from self-employment Social Security taxes under 26 U.S.C. \S~1402(h) (\href{https://supreme.justia.com/cases/federal/us/406/205/}{406 U.S. at 222 \& n.11}). They paralleled Thomas Jefferson's ideal of the ``sturdy yeoman,'' and compulsory schooling past the eighth grade was itself a relatively recent development less than 60 years old, historically linked to child labor laws whose concerns did not apply to children working under parental supervision on family farms (\href{https://supreme.justia.com/cases/federal/us/406/205/}{406 U.S. at 225--29}). Wisconsin's fear that Amish youth who later left the community would become societal burdens was ``highly speculative'' and unsupported by record evidence (\href{https://supreme.justia.com/cases/federal/us/406/205/}{406 U.S. at 224--25}).

\subsection{Part IV: Rejection of Broad \emph{Parens Patriae} Authority and Non-Resolution of Parent--Child Conflicts}

Wisconsin invoked \href{https://supreme.justia.com/cases/federal/us/321/158/}{\emph{Prince v. Massachusetts}}, 321 U.S. 158 (1944), to argue that the State as \emph{parens patriae} could secure secondary education for children regardless of parental wishes. The Court confined \emph{Prince} to conduct posing a ``substantial threat to public safety, peace or order'' or jeopardizing a child's physical or mental health---conditions absent here (\href{https://supreme.justia.com/cases/federal/us/406/205/}{406 U.S. at 229--30, 233--34}). Crucially, the Court stressed that \href{https://supreme.justia.com/cases/federal/us/268/510/}{\emph{Pierce}} ``stands as a charter of the rights of parents to direct the religious upbringing of their children,'' and that ``when the interests of parenthood are combined with a free exercise claim of the nature revealed by this record, more than merely a `reasonable relation to some purpose within the competency of the State' is required'' (\href{https://supreme.justia.com/cases/federal/us/406/205/}{406 U.S. at 233}). Responding to Justice Douglas's partial dissent, the majority emphasized that the parents---not the children---were being criminally prosecuted, and that because no child in the record had expressed a desire to attend high school contrary to parental wishes, the Court neither reached nor decided how a genuine parent--child conflict should be resolved (\href{https://supreme.justia.com/cases/federal/us/406/205/}{406 U.S. at 230--32}).

\subsection{Part V: Establishment Clause Accommodation and Narrowing Caution}

In footnote~22, the Court rejected any suggestion that exempting the Amish violated the Establishment Clause, explaining that accommodating a centuries-old sect's survival reflects governmental neutrality rather than sponsorship or active involvement (\href{https://supreme.justia.com/cases/federal/us/406/205/}{406 U.S. at 234 n.22}). Finally, cautioning that ``courts are not school boards or legislatures,'' the Court stressed that the Amish had made a ``convincing showing, one that probably few other religious groups or sects could make,'' and left room for Wisconsin to adopt reasonable vocational-education standards consistent with Free Exercise (\href{https://supreme.justia.com/cases/federal/us/406/205/}{406 U.S. at 235--36}).

\subsection{Concurring Opinion (Justice Stewart)}

Justice Stewart wrote briefly to underscore that the case presented only the constitutionality of branding Amish parents as criminals, and ``in no way involves any questions regarding the right of the children of Amish parents to attend public high schools'' (\href{https://supreme.justia.com/cases/federal/us/406/205/}{406 U.S. at 237}). Quoting Frieda Yoder's cross-examination confirming that religion was her sole reason for not attending high school, he concluded that the record did not raise the parent--child conflict issue pressed by Justice Douglas (\href{https://supreme.justia.com/cases/federal/us/406/205/}{406 U.S. at 237}).

\subsection{Concurring Opinion (Justice White)}

Justice White described the case as ``close'' and stressed that it would be ``a very different case'' if a religion forbade all formal schooling (\href{https://supreme.justia.com/cases/federal/us/406/205/}{406 U.S. at 238, 240}). Quoting \emph{Brown v. Board of Education}, 347 U.S. 483, 493 (1954), he highlighted the State's legitimate interest in developing children's latent talents---noting that some Amish youth leave the faith and might aspire to become ``nuclear physicists, ballet dancers, computer programmers, or historians'' (\href{https://supreme.justia.com/cases/federal/us/406/205/}{406 U.S. at 238--40}). He joined the Court because the Amish deviation from Wisconsin's age-16 minimum was ``relatively slight'' (at most two years beyond eighth grade), the sincerity and severity of the burden were undisputed, and the State had not shown that leaving school after eighth grade would intellectually stultify the children (\href{https://supreme.justia.com/cases/federal/us/406/205/}{406 U.S. at 238, 240--41}).

\subsection{Opinion Dissenting in Part (Justice Douglas)}

Justice Douglas joined the judgment as to Jonas Yoder because Frieda Yoder had testified to her own religious objections, but dissented as to Adin Yutzy and Wallace Miller because neither Vernon Yutzy nor Barbara Miller had been heard (\href{https://supreme.justia.com/cases/federal/us/406/205/}{406 U.S. at 243, 246}). He advanced three major objections:
\begin{enumerate}
  \item \textbf{Children's Independent Constitutional Personhood:} Citing \emph{In re Gault}, 387 U.S. 1 (1967), \emph{Tinker v. Des Moines Independent Community School District}, 393 U.S. 503 (1969), and \emph{West Virginia State Board of Education v. Barnette}, 319 U.S. 624 (1943), Douglas argued that 14- and 15-year-old students are ``persons'' under the Bill of Rights with independent religious and educational interests (\href{https://supreme.justia.com/cases/federal/us/406/205/}{406 U.S. at 243--45}). Because ``[i]t is the future of the student, not the future of the parents, that is imperiled''---and a child ``may want to be a pianist or an astronaut or an oceanographer''---courts should canvass the mature child's views before granting an exemption (\href{https://supreme.justia.com/cases/federal/us/406/205/}{406 U.S. at 244--46}).
  \item \textbf{Irrelevance of the Amish ``Law and Order'' Record:} Douglas criticized the majority's reliance on the Amish community's low crime rate and idyllic agrarianism as constitutionally irrelevant and factually romanticized, noting that ``[a] religion is a religion irrespective of what the misdemeanor or felony records of its members might be'' (\href{https://supreme.justia.com/cases/federal/us/406/205/}{406 U.S. at 246--47 \& n.5}).
  \item \textbf{Exclusion of Secular/Philosophical Conscience (Thoreau):} Douglas faulted the majority's dismissal of Henry David Thoreau's philosophical convictions as inconsistent with the broad, functional definition of ``religious'' conscience adopted in conscientious-objector cases such as \emph{United States v. Seeger}, 380 U.S. 163 (1965), and \href{https://supreme.justia.com/cases/federal/us/398/333/}{\emph{Welsh v. United States}}, 398 U.S. 333 (1970) (\href{https://supreme.justia.com/cases/federal/us/406/205/}{406 U.S. at 247--49}).
\end{enumerate}

\subsection{Subsequent Developments}

\begin{itemize}

\item \textbf{Early Narrowing:} Over the next two decades, however, the Court narrowed its reach: \href{https://supreme.justia.com/cases/federal/us/455/252/}{\emph{United States v. Lee}}, 455 U.S. 252, 258--61 (1982), denied an Amish employer a Social Security tax exemption, and \href{https://supreme.justia.com/cases/federal/us/485/439/}{\emph{Lyng v. Northwest Indian Cemetery Protective Ass'n}}, 485 U.S. 439, 456--57 (1988), confined \emph{Yoder} to direct governmental coercion rather than internal government operations.

\item \textbf{\emph{Smith} and the ``Hybrid Situation'':} In \href{https://supreme.justia.com/cases/federal/us/494/872/}{\emph{Employment Division v. Smith}}, 494 U.S. 872, 878--85 (1990), the Court largely repudiated \emph{Sherbert}--\emph{Yoder} strict scrutiny for neutral, generally applicable laws, but preserved \emph{Yoder} as a ``hybrid situation'' combining Free Exercise with the parental right to direct education recognized in \href{https://supreme.justia.com/cases/federal/us/268/510/}{\emph{Pierce v. Society of Sisters}}, 268 U.S. 510 (1925) (494 U.S. at 881--82 \& n.1).

\item \textbf{Revitalization in \emph{Mahmoud v. Taylor}:} In \href{https://www.law.cornell.edu/supremecourt/text/24-297}{\emph{Mahmoud v. Taylor}}, 606 U.S. 522 (2025) (6--3, Alito, J.), the Court decisively revitalized \emph{Yoder}, rejecting lower-court views that it was ``\emph{sui generis}'' and holding instead that it ``embodies a robust principle of general applicability.'' Granting a preliminary injunction to Muslim, Catholic, and Ukrainian Orthodox parents seeking advance notice and opt-outs from K--5 ``LGBTQ+-inclusive'' storybooks, the Court held that normative instruction posing ``a very real threat of undermining'' parental religious formation imposes a burden ``of the exact same character as that in \emph{Yoder},'' triggering strict scrutiny \emph{regardless} of \emph{Smith} neutrality and without deciding ``hybrid rights'' (606 U.S. 522, n.14). Justice Thomas concurred on \emph{Yoder}'s historical-pedigree inquiry, while Justice Sotomayor dissented (joined by Justices Kagan and Jackson), warning that treating ``mere exposure'' as coercion contradicts \emph{Lyng} and risks administrative chaos.

\end{itemize}

\section{Examples: Future Applications}
\begin{enumerate}

\item \textbf{Same-Side (Exemption Required): Raising the Attendance Age to 18.} A State raises its compulsory-attendance age from 16 to 18 and prosecutes Old Order Amish parents whose 16- and 17-year-olds are learning farming and carpentry at home under parental supervision. The burden is the one the Court found in \textit{Yoder}: formal high school removes adolescents from the community during the formative years before adult baptism. The State's interests in citizenship and self-sufficiency are already served by eight years of schooling plus the Amish vocational system, and it cannot show with particularity how two more years of classroom attendance for this small group would harm those interests. Under \textit{Yoder}'s ``interests of the highest order'' standard, the prosecutions fail (\href{https://supreme.justia.com/cases/federal/us/406/205/}{Justia}).

\item \textbf{Same-Side (Strict Scrutiny Triggered): Normative Classroom Instruction.} Devout parents seek advance notice and the chance to excuse their children from a mandatory fourth-grade unit with teacher scripts correcting traditional gender-role views. The parents are not challenging what the school teaches other children; they seek to keep their own children out of instruction designed to change the religious views the parents are trying to instill. After \textit{Mahmoud v. Taylor} read \textit{Yoder} as a general principle rather than an Amish-specific exception, instruction of this kind is treated as a burden on parents' right to direct their children's religious upbringing that triggers strict scrutiny even though the curriculum applies to everyone. The school district must show a compelling reason for refusing an opt-out.

\item \textbf{Opposite-Side (No Exemption): Secular Communitarians.} Parents in a self-sufficient, law-abiding, non-theistic agrarian commune inspired by Thoreau withdraw their 15-year-olds after eighth grade for structured farm-vocational training. Under \emph{Yoder}, their claim \emph{fails} strict scrutiny because ``purely secular considerations'' do not trigger the Religion Clauses (\href{https://supreme.justia.com/cases/federal/us/406/205/}{406 U.S. at 215--16}), leaving only rational-basis review under \emph{Pierce}---though Douglas's dissent (\href{https://supreme.justia.com/cases/federal/us/398/333/}{\emph{Welsh}}) and Establishment Clause parity provide sharp counterarguments.

\item \textbf{Opposite-Side (No Exemption): A Faith That Forbids All Schooling.} Members of a small religious community refuse to send their children to any school, or to provide any home instruction, because their faith forbids formal learning of any kind; their children reach adolescence unable to read. The State prosecutes the parents under its compulsory-attendance law. \textit{Yoder} rested heavily on the Amish acceptance of eight years of elementary schooling and on the ``relatively slight'' deviation at issue, and Justice White's concurrence, joined by two members of the majority, stressed that it would be ``a very different case'' if a religion forbade all formal schooling. The State's interest in basic literacy and the child's future is of the highest order and cannot be served by a less restrictive means, so the prosecutions likely succeed.

\item \textbf{Fence-Sitter (Uncertain): A Vaccine Requirement With Medical Exemptions.} Devout parents seek a religious exemption from a mandatory measles-vaccine rule permitting medical exemptionsas a condition of school attendance. The State can rely on \textit{Prince v. Massachusetts}, which \textit{Yoder} distinguished as involving a threat to public health and safety, and argue that a communicable-disease rule does not interfere with religious formation the way high school did for the Amish. But the parents can argue that allowing medical exemptions while refusing religious ones means the rule is not generally applicable under \textit{Fulton v. City of Philadelphia}, 593 U.S. 522 (2021), which would trigger strict scrutiny regardless of \textit{Yoder}. Whether the State's interest in preventing outbreaks survives that scrutiny is genuinely contested.

\end{enumerate}

\section{Critique}
\textbf{``Good Citizen'' Bias:} Relying on the Amish's 300-year pedigree and low crime rate (\href{https://supreme.justia.com/cases/federal/us/406/205/}{406 U.S. at 222, 235--36}) risks disfavoring newer or unpopular faiths and entangling courts in evaluating religious lifestyles.

\textbf{The Thoreau Boundary:} Excluding secular conscience (\href{https://supreme.justia.com/cases/federal/us/406/205/}{406 U.S. at 215--16}) sharpens tension between Free Exercise accommodation and Establishment Clause neutrality.

\textbf{Children's Autonomy and Post-\emph{Mahmoud} Scope:} Deferring children's independent interests (Douglas's dissent) and detaching \emph{Yoder} from existential sect survival in \emph{Mahmoud} raise acute limiting-principle questions for public-school curricula.

\section{Key Quotations}
\begin{quote}
``The essence of all that has been said and written on the subject is that only those interests of the highest order and those not otherwise served can overbalance legitimate claims to the free exercise of religion.''
\end{quote}
\noindent --- Burger, C.J., \href{https://supreme.justia.com/cases/federal/us/406/205/}{406 U.S. at 215}; \href{https://www.law.cornell.edu/supremecourt/text/406/205}{Cornell LII}.

\begin{quote}
``A way of life, however virtuous and admirable, may not be interposed as a barrier to reasonable state regulation of education if it is based on purely secular considerations; to have the protection of the Religion Clauses, the claims must be rooted in religious belief.''
\end{quote}
\noindent --- Burger, C.J., \href{https://supreme.justia.com/cases/federal/us/406/205/}{406 U.S. at 215}; \href{https://www.law.cornell.edu/supremecourt/text/406/205}{Cornell LII}.

\begin{quote}
``As the record shows, compulsory school attendance to age 16 for Amish children carries with it a very real threat of undermining the Amish community and religious practice as they exist today; they must either abandon belief and be assimilated into society at large or be forced to migrate to some other and more tolerant region.''
\end{quote}
\noindent --- Burger, C.J., \href{https://supreme.justia.com/cases/federal/us/406/205/}{406 U.S. at 218}; \href{https://www.law.cornell.edu/supremecourt/text/406/205}{Cornell LII}.

\begin{quote}
``A regulation neutral on its face may, in its application, nonetheless offend the constitutional requirement for governmental neutrality if it unduly burdens the free exercise of religion.''
\end{quote}
\noindent --- Burger, C.J., \href{https://supreme.justia.com/cases/federal/us/406/205/}{406 U.S. at 220}; \href{https://www.law.cornell.edu/supremecourt/text/406/205}{Cornell LII}.

\begin{quote}
``It is the future of the student, not the future of the parents, that is imperiled by today's decision. If a parent keeps his child out of school beyond the grade school, then the child will be forever barred from entry into the new and amazing world of diversity that we have today.''
\end{quote}
\noindent --- Douglas, J., dissenting in part, \href{https://supreme.justia.com/cases/federal/us/406/205/}{406 U.S. at 245}; \href{https://www.law.cornell.edu/supremecourt/text/406/205}{Cornell LII}.


\end{document}
